%==========================================================================
%  TRES0312 Fysikk — Foiler-mal
%
%  MØrk/lys modus: kommenter ut én av de to linjene nedenfor.
%==========================================================================
\newif\ifdarkmode
\darkmodetrue      % ← mørk modus
%\darkmodefalse   % ← lys modus  (fjern % foran og legg % foran linjen over)

\documentclass[aspectratio=169, 12pt]{beamer}

\usepackage{amd-beamer}
\usepackage[european]{circuitikz}
\usetikzlibrary{overlay-beamer-styles}
\usetikzlibrary{calc}

%--------------------------------------------------------------------------
%  DOKUMENTINFO
%--------------------------------------------------------------------------
\title{Ohms lov, koblinger og effekt}
\subtitle{TRES0312 --- kapittel 9.3--9.6}
\author{Ben David Normann}
\date{\today}

%--------------------------------------------------------------------------
%  SEKSJONSKONTEKST
%--------------------------------------------------------------------------
\setcounter{section}{9}   % Kapittel 9: "Elektrisitet"

\begin{document}

%--------------------------------------------------------------------------
%  SLIDE 1 — TITTELSIDE
%--------------------------------------------------------------------------
\begin{frame}[plain]
  \titlepage
  \dekkerdelkap{kapittel 9.3--9.6}
  \vspace{0.3cm}
  \begin{center}
  {\bfseries\color{repcol} Første time: Ohms lov, serie- og parallellkobling. Etter pausen: Kirchhoffs lover og effekt.}
  \end{center}
\end{frame}

%--------------------------------------------------------------------------
%  SLIDE 2 — Tre kjappe fra forrige gang
%--------------------------------------------------------------------------
\begin{frame}{}

  \repetisjon{Tre kjappe fra forrige gang}{%
    \begin{enumerate}[1.]
      \item Hvor mange elektroner må til for å få ladningen $1\,\mathrm{C}$?
      \item Hva er strøm, og hva er spenning?
      \item Hvordan kobler vi et amperemeter og et voltmeter?
    \end{enumerate}
  }

\end{frame}

%==========================================================================
%  DEL 1: OHMS LOV OG KOBLINGER (9.3 OG 9.4)
%==========================================================================

%--------------------------------------------------------------------------
%  SLIDE 3 — Til ettertanke: ledning og glødetråd
%--------------------------------------------------------------------------
\begin{frame}{Resistans}

  \tanke{}{%
  Hvorfor blir ledningen til en vannkoker nesten ikke varm, mens glødetråden i en gammeldags lyspære blir så varm at den lyser?
  }
  \pause
  \svar{
  Glødetråden har mye større \textit{resistans} enn ledningen. Elektronene kolliderer med atomene og gir fra seg energi. Jo større resistans, desto mer energi blir til varme der.
  }

\end{frame}

%--------------------------------------------------------------------------
%  SLIDE 4 — Definisjon: Ohms lov
%--------------------------------------------------------------------------
\begin{frame}{Ohms lov}

  \defnat{9.3.1}{Ohms lov}{
  \[
    U = R\cdot I \qquad [R] = \Omega = \mathrm{V/A}
  \]
  der $U$ er spenningen over komponenten, $I$ er strømmen gjennom den og $R$ er resistansen.
  }
  \pause
  \begin{itemize}
    \item Tre former: $U = RI$, \ $I = U/R$, \ $R = U/I$.
    \item Stor resistans gir liten strøm for samme spenning.
  \end{itemize}
  \pause
  \merk{Kortslutning}{%
  Med nesten null resistans blir strømmen $I = U/R$ svært stor. Det er derfor vi har sikringer.
  }

\end{frame}

%--------------------------------------------------------------------------
%  SLIDE 5 — Aktivitet: Ohms lov
%--------------------------------------------------------------------------
%  Fasit: a) I = 12/33 = 0.364 A ≈ 0.36 A.
%         b) R = 230/0.43 = 535 Ω ≈ 5.3e2 Ω.
\begin{frame}{Aktivitet}

  \aktivitet{Ohms lov}{%
  \begin{enumerate}[a)]
    \item En motstand på $33\,\Omega$ kobles til et batteri på $12\,\mathrm{V}$. Hvor stor blir strømmen?
    \item En lyspære kobles til $230\,\mathrm{V}$, og strømmen er $0.43\,\mathrm{A}$. Hva er resistansen i pæra?
  \end{enumerate}
  }

\end{frame}

%--------------------------------------------------------------------------
%  SLIDE 6 — Til ettertanke: julelys
%--------------------------------------------------------------------------
\begin{frame}{Serie- og parallellkobling}

  \tanke{}{%
  I et gammeldags julelys slukker hele lenka når én pære går. Men når én lampe i stua går, lyser de andre fortsatt. Hva er forskjellen?
  }
  \pause
  \svar{
  Julelyset er koblet i \textit{serie}: strømmen har bare én vei, og et brudd åpner hele kretsen. Lampene i huset er koblet i \textit{parallell}: hver lampe har sin egen vei.
  }

\end{frame}

%--------------------------------------------------------------------------
%  SLIDE 7 — Serie og parallell
%--------------------------------------------------------------------------
\begin{frame}{Serie- og parallellkobling}

  \begin{columns}[T]
    \begin{column}{0.48\textwidth}
      \centering
      \textbf{Serie}\\[2mm]
      \begin{circuitikz}[color=textcolor, scale=0.75, transform shape]
        \draw (0,0) to[battery1, l=$U$] (0,2.5)
              to[R, l=$R_1$] (2.5,2.5)
              to[R, l=$R_2$] (2.5,0) -- (0,0);
      \end{circuitikz}
      \begin{itemize}
        \item Samme \textbf{strøm} gjennom alle
        \item Spenningen deles: $U = U_1 + U_2$
      \end{itemize}
    \end{column}
    \begin{column}{0.48\textwidth}
      \centering
      \textbf{Parallell}\\[2mm]
      \begin{circuitikz}[color=textcolor, scale=0.75, transform shape]
        \draw (0,0) to[battery1, l=$U$] (0,2.5) -- (3.5,2.5);
        \draw (1.8,2.5) to[R, l=$R_1$] (1.8,0);
        \draw (3.5,2.5) to[R, l=$R_2$] (3.5,0) -- (0,0);
      \end{circuitikz}
      \begin{itemize}
        \item Samme \textbf{spenning} over alle
        \item Strømmen deles: $I = I_1 + I_2$
      \end{itemize}
    \end{column}
  \end{columns}

\end{frame}

%--------------------------------------------------------------------------
%  SLIDE 8 — Definisjon: serie og parallell
%--------------------------------------------------------------------------
\begin{frame}{Serie- og parallellkobling}

  \defnat{9.4.1}{Serie- og parallellkobling av motstander}{
  \[
    \text{Serie: } R_\text{tot} = R_1 + R_2 + \dots \qquad\quad
    \text{Parallell: } \frac{1}{R_\text{tot}} = \frac{1}{R_1} + \frac{1}{R_2} + \dots
  \]
  For to motstander i parallell: $R_\text{tot} = \dfrac{R_1R_2}{R_1 + R_2}$.
  }
  \pause
  \merk{Husk:}{%
  Parallellkobling gir strømmen flere veier. $R_\text{tot}$ blir alltid \textit{mindre} enn den minste motstanden.
  }

\end{frame}

%--------------------------------------------------------------------------
%  SLIDE 9 — Aktivitet: spenningsdeling i serie
%--------------------------------------------------------------------------
%  Fasit: R_tot = 9.0 Ω. I = 18/9.0 = 2.0 A. U1 = 3.0*2.0 = 6.0 V, U2 = 6.0*2.0 = 12 V.
\begin{frame}{Aktivitet}

  \aktivitet{Spenningsdeling i serie}{%
  To motstander, $R_1 = 3.0\,\Omega$ og $R_2 = 6.0\,\Omega$, er koblet i serie til et batteri på $18\,\mathrm{V}$. Finn strømmen og spenningen over hver motstand. Kontroller at spenningene til sammen blir $18\,\mathrm{V}$.
  }

\end{frame}

%--------------------------------------------------------------------------
%  SLIDE 10 — Aktivitet: motstander i parallell
%--------------------------------------------------------------------------
%  Fasit: R_tot = 4.0*12/16 = 3.0 Ω. I1 = 12/4.0 = 3.0 A, I2 = 12/12 = 1.0 A, I = 4.0 A.
%         Kontroll: 12 V/3.0 Ω = 4.0 A.
\begin{frame}{Aktivitet}

  \aktivitet{Motstander i parallell}{%
  To motstander, $R_1 = 4.0\,\Omega$ og $R_2 = 12\,\Omega$, er koblet i parallell til et batteri på $12\,\mathrm{V}$.
  \begin{enumerate}[a)]
    \item Finn den totale resistansen.
    \item Finn strømmen gjennom hver motstand og strømmen fra batteriet.
  \end{enumerate}
  }

\end{frame}

%--------------------------------------------------------------------------
%  SLIDE 11 — Oppsummering: del 1
%--------------------------------------------------------------------------
\begin{frame}{}

  \oppsum{Ohms lov og koblinger}{%
  \begin{enumerate}[1.]
    \item[] Ohms lov: $U = RI$.
    \item[] Serie: samme strøm, $R_\text{tot} = R_1 + R_2 + \dots$
    \item[] Parallell: samme spenning, $1/R_\text{tot} = 1/R_1 + 1/R_2 + \dots$
  \end{enumerate}
  }

\end{frame}

%--------------------------------------------------------------------------
%  SLIDE 12 — PAUSE
%--------------------------------------------------------------------------
\begin{frame}
\begin{center}
  \Huge{Pause}\\[8pt]
  {\large Etter pausen: Kirchhoffs lover og effekt}
\end{center}
\end{frame}

%==========================================================================
%  DEL 2: KIRCHHOFFS LOVER OG EFFEKT (9.5 OG 9.6)
%==========================================================================

%--------------------------------------------------------------------------
%  SLIDE 13 — Definisjon: Kirchhoffs lover
%--------------------------------------------------------------------------
\begin{frame}{Kirchhoffs lover}

  \defnat{9.5.1}{Kirchhoffs lover}{
  \textbf{Strømloven:} I et forgreningspunkt er $\sum I_\text{inn} = \sum I_\text{ut}$.\\[2mm]
  \textbf{Spenningsloven:} Rundt en lukket sløyfe er $\sum U_\text{kilder} = \sum U_\text{fall}$.
  }
  \pause
  \begin{itemize}
    \item Strømloven: \textbf{ladning er bevart}.
    \item Spenningsloven: \textbf{energi er bevart}. En runde i sløyfa tar deg tilbake til samme «høyde».
  \end{itemize}

\end{frame}

%--------------------------------------------------------------------------
%  SLIDE 14 — Slik regner du på en sammensatt krets
%--------------------------------------------------------------------------
\begin{frame}{Sammensatte kretser}

  \begin{enumerate}[1.]
    \item<1-> Tegn kretsen, og marker strømmene.
    \item<2-> Finn ut hva som er i serie og hva som er i parallell.
    \item<3-> Forenkle trinn for trinn til én total resistans.
    \item<4-> Finn strømmen fra batteriet: $I = U/R_\text{tot}$.
    \item<5-> Gå baklengs: spenningen over hver del, og så strømmene i grenene.
  \end{enumerate}

\end{frame}

%--------------------------------------------------------------------------
%  SLIDE 15 — Aktivitet: krets med serie og parallell
%--------------------------------------------------------------------------
%  Fasit: R_p = 10*15/25 = 6.0 Ω. R_tot = 4.0 + 6.0 = 10 Ω.
%         I0 = 20/10 = 2.0 A. U_p = 6.0*2.0 = 12 V. U1 = 4.0*2.0 = 8.0 V (8 + 12 = 20 V).
%         I2 = 12/10 = 1.2 A, I3 = 12/15 = 0.80 A. Sum 2.0 A.
\begin{frame}{Aktivitet}

  \begin{columns}[c]
    \begin{column}{0.55\textwidth}
      \aktivitet{Krets med serie og parallell}{%
      \small
      Batteriet har polspenningen $U_0 = 20\,\mathrm{V}$. $R_1 = 4.0\,\Omega$ er i serie med en parallellkobling av $R_2 = 10\,\Omega$ og $R_3 = 15\,\Omega$.
      \begin{enumerate}[a)]
        \item Finn $R_\text{tot}$ og strømmen $I_0$.
        \item Finn spenningen over parallellkoblingen.
        \item Finn $I_2$ og $I_3$. Kontroller med strømloven.
      \end{enumerate}
      }
    \end{column}
    \begin{column}{0.43\textwidth}
      \centering
      \begin{circuitikz}[color=textcolor, scale=0.75, transform shape]
        \draw (0,0) to[battery1, l=$U_0$] (0,3)
              to[R, l=$R_1$, i>^=$I_0$] (3,3) -- (5.5,3);
        \draw (3,3) to[R, l_=$R_2$, i>_=$I_2$] (3,0);
        \draw (5.5,3) to[R, l=$R_3$, i>^=$I_3$] (5.5,0);
        \draw (5.5,0) -- (0,0);
        \fill (3,3) circle (2pt) (3,0) circle (2pt);
      \end{circuitikz}
    \end{column}
  \end{columns}

\end{frame}

%--------------------------------------------------------------------------
%  SLIDE 16 — Til ettertanke: vannkoker og lyspære
%--------------------------------------------------------------------------
\begin{frame}{Elektrisk effekt}

  \tanke{}{%
  På en vannkoker står det $2000\,\mathrm{W}$. På en lyspære står det $60\,\mathrm{W}$. Hva forteller disse tallene?
  }
  \pause
  \svar{
  Hvor mye energi apparatet bruker per sekund. Vannkokeren bruker $2000\,\mathrm{J}$ hvert sekund, lyspæra $60\,\mathrm{J}$.
  }

\end{frame}

%--------------------------------------------------------------------------
%  SLIDE 17 — Definisjon: elektrisk effekt
%--------------------------------------------------------------------------
\begin{frame}{Elektrisk effekt}

  \defnat{9.6.1}{Elektrisk effekt}{
  \[
    P = U\cdot I \qquad [P] = \mathrm{W} = \mathrm{J/s}
  \]
  Med Ohms lov: $P = RI^2 = \dfrac{U^2}{R}$.
  }
  \pause
  Energi er effekt ganger tid, $E = P\cdot t$. På strømregningen måles den i kilowattimer:
  \graabox{$\displaystyle 1\,\mathrm{kWh} = 1000\,\mathrm{W}\cdot 3600\,\mathrm{s} = 3.6\cdot 10^6\,\mathrm{J}$}

\end{frame}

%--------------------------------------------------------------------------
%  SLIDE 18 — Aktivitet: vannkoker
%--------------------------------------------------------------------------
%  Fasit: a) I = 2200/230 = 9.57 A ≈ 9.6 A.
%         b) R = 230^2/2200 = 24.0 Ω ≈ 24 Ω.
%         c) E = 2.2 kW * (4.0/60) h = 0.147 kWh ≈ 0.15 kWh = 5.3e5 J.
\begin{frame}{Aktivitet}

  \aktivitet{Vannkoker}{%
  En vannkoker har effekten $2200\,\mathrm{W}$ og er koblet til $230\,\mathrm{V}$.
  \begin{enumerate}[a)]
    \item Hvor stor er strømmen gjennom vannkokeren?
    \item Hva er resistansen i varmeelementet?
    \item Den står på i $4.0\,\mathrm{min}$. Hvor mye energi bruker den, i kWh og i joule?
  \end{enumerate}
  }

\end{frame}

%--------------------------------------------------------------------------
%  SLIDE 19 — Aktivitet: hvilken pære lyser sterkest?
%--------------------------------------------------------------------------
%  Fasit: Serie: R_tot = 8.0 Ω, I = 12/8.0 = 1.5 A, P = 4.0*1.5^2 = 9.0 W per pære.
%         Parallell: P = 12^2/4.0 = 36 W per pære. Sterkest i parallell (4 ganger).
\begin{frame}{Aktivitet}

  \aktivitet{Hvilken pære lyser sterkest?}{%
  To like lyspærer med resistans $4.0\,\Omega$ kobles til et batteri på $12\,\mathrm{V}$, først i serie og så i parallell. Finn effekten i hver pære i de to tilfellene. Når lyser pærene sterkest?
  }
  \pause
  \merk{Derfor:}{%
  Lampene i et hus er koblet i parallell. Da får hver lampe hele spenningen fra stikkontakten.
  }

\end{frame}

%--------------------------------------------------------------------------
%  SLIDE 20 — Oppsummering: del 2
%--------------------------------------------------------------------------
\begin{frame}{}

  \oppsum{Kirchhoffs lover og effekt}{%
  \begin{enumerate}[1.]
    \item[] Strømloven: $\sum I_\text{inn} = \sum I_\text{ut}$. Spenningsloven: $\sum U_\text{kilder} = \sum U_\text{fall}$.
    \item[] Effekt: $P = UI = RI^2 = U^2/R$.
    \item[] Energi: $E = Pt$, og $1\,\mathrm{kWh} = 3.6\cdot 10^6\,\mathrm{J}$.
  \end{enumerate}
  }

\end{frame}

%--------------------------------------------------------------------------
%  SLIDE 21 — Neste gang
%--------------------------------------------------------------------------
\begin{frame}
\begin{center}
  \Huge{Neste gang: repetisjon og eksamensoppgaver}
  \vspace{6mm}
\end{center}
\end{frame}

\end{document}
